\section{Introduction}
\label{sec:intro}

%% 结构（见 outline_v1.md §1）：¶1 问题重要性 + 目标是什么 → ¶2 rs-fMRI 在组水平产出了什么
%% → ¶3 枢纽：组水平差异 ≠ 个体水平预测，而后者塌了 → ¶4 两个已确立的原因 + 并非无解
%% → ¶5 方法登场 + 自动化搜索带来的危险 → ¶6 数据集 + contributions + roadmap
%%
%% 2026-08-09 修订（v2.1）：删除原 ¶5「第三个原因：算法逐篇手工设计、评估严格度决定数字」。
%%   理由（用户判断）：该段以「各篇论文在自选划分/子样本/流水线上评估」作为指责，而本文自己
%%   同样在自选的划分上评估，论证会反噬；且评估严格度这条线 §2.1 已完整承载，§1 重复无收益。
%%   ¶4 与 ¶6 直接承接：由「方法与评估都做对时是有空间的」→「但方法一直是一次一个架构地设计」，
%%   一句带过，不再展开严格度论证，也不再前指 §2.1。
%% 纪律：不作优先权宣称（no "first"/"the first to"）；不写协议机制细节（留 §5）。
%%
%% 2026-08-09 全面重写（v2）。依据 reference/mdd-related 下 13 篇论文 Introduction 的调研：
%%   (a) 这一体裁无一篇用短促独立句作修辞，均为多从句匀速学术散文——v1 的碎句已全部并句；
%%   (b) 大刊（Gallo/Mol Psych、Zheng/PLOS Digit Health）的 intro 极简、只给聚合量级，
%%       把逐研究数字留给正文，v1 的 ¶3 数字堆叠已按此压缩；
%%   (c) Yamashita 2020 自陈：临床上区分 MDD 与健康人对有经验的医生并不难，marker 的价值
%%       在于机制与治疗选择——¶1 末句正面回应这一点，避免「用分类器替代访谈」的弱论证。
%% 新增论证（v1 没有的）：组水平一致性与个体水平可判别性是两个量，且组水平连符号都未settled
%%   （Kaiser 元分析报 DMN 高连接，Yan 在最大队列报复发患者 DMN 低连接）——这是 ¶2→¶3 的枢纽。
%% 新增引用：kaiser2015（JAMA Psychiatry 元分析，Crossref + Europe PMC 双重核实，方向性结论
%%   取自其 Results 原文：frontoparietal 内 hypoconnectivity、default network 内 hyperconnectivity）。

Depressive disorders ranked second among all causes of years lived with disability in the Global
Burden of Disease Study 2019 \citep{gbd2019mental}. The diagnosis rests on clinical interview and
symptom-scale scores; no laboratory or imaging measurement enters it. The category so defined is
broad and is applied with limited reproducibility: in the DSM-5 field trials, two clinicians who
independently assessed the same patient agreed on a diagnosis of major depressive disorder at a
pooled intraclass $\kappa$ of 0.28 \citep{regier2013dsm5}. The case for a biological
marker therefore rests less on any inability of clinicians to recognize depression than on what a
measurement independent of the interview would supply: a quantity that can be related to mechanism,
to treatment response, and to the boundaries of a category that symptom criteria have not sharpened.

Resting-state functional MRI has been the most actively pursued candidate for such a measurement:
it records spontaneous coordination between brain regions without requiring task performance, and
can therefore be acquired under uniform conditions across patients whose ability to comply with a
task would differ. Group comparisons have produced a broadly consistent account of which systems
are implicated. A meta-analysis pooling 27 data sets associated the disorder with hypoconnectivity
within the frontoparietal control network and with altered coupling between that network and the
default mode network \citep{kaiser2015}; analyses of the REST-meta-MDD consortium reported reduced
default mode connectivity in patients with recurrent illness \citep{yan2019pnas}; and thalamic
hyperconnectivity was the dominant discriminative feature in each of two independently collected
consortia \citep{gallo2023}. The direction of the difference is less settled than its location: the
meta-analytic estimate for the default mode network is one of hyperconnectivity, while the largest
single consortium reports the opposite in its recurrent patients.

That the groups differ on average does not establish that an individual scan can be assigned to one
of them, and it is the latter that any use of such a marker would require. Measured directly,
individual-level discrimination has not held up. A meta-analysis of resting-state classification
studies reported 85\% sensitivity and 83\% specificity across samples averaging roughly 28 patients
each \citep{kambeitz2017}. Posed on the two largest available multi-site consortia, the same problem
yields 61\% balanced accuracy, and falls to 53--56\% when one consortium is used for training and
the other for testing \citep{gallo2023}. Stratifying studies by sample size shows the same relation
from the other side, with the larger studies reporting the lower sensitivities and specificities
\citep{chen2023meta}. Accuracy that falls as samples grow is the reverse of what a learning curve
predicts, and is what a literature reports when each published entry is a maximum taken over choices
that go unrecorded \citep{schnack2016, marek2022}.

Two reasons for this are established. Depression is a syndrome rather than a lesion, so patients who
share a diagnosis need not share a neurobiology, and what a classifier can learn is bounded by the
consistency of the labels it is trained on. Differences between acquisition sites are the second
reason, and they have been measured rather than asserted: using participants who travelled between
scanners, \citet{yamashita2019} decomposed site differences into a sampling component and a
measurement component, and found the effect of both on pairwise functional connectivity to be
greater than or equal to the effect of the psychiatric disorders themselves. Neither reason makes
the problem intractable when the method and the evaluation are designed for it. After explicit
harmonization, a marker trained on four sites transferred to 521 participants from five independent
sites at approximately 70\% accuracy \citep{yamashita2020}.

What that method should be, however, has been approached one architecture at a time, each proposal
a single point in a large space of candidate algorithms that has been sampled sparsely. We take a
different route: an LLM-driven evolutionary coding agent writes, executes, and revises complete
classification programs, so that the choice of connectome view, the construction of features, the
network architecture, the loss, and the training schedule are decided by the search rather than
fixed in advance. Automating the search also automates the analytic freedom that produced the
pattern described above. Where a human explores a few dozen variants, an agent explores hundreds,
in whichever direction raises the reported number. The evaluation must therefore be hardened before
the search begins, not audited after it ends.

The work is carried out on REST-meta-MDD / DIRECT Phase II, the largest openly shared multi-site
resting-state fMRI resource for depression: 1{,}660 patients with major depressive disorder and
1{,}341 healthy controls from 23 cohorts, 3{,}001 participants in total, with every site processed
through a common pipeline \citep{direct2}.

\paragraph{Contributions.}
\begin{itemize}
  \item \textbf{An end-to-end algorithm discovery loop for connectome classification.} The agent
    designs executable programs rather than tuning the hyperparameters of a fixed model, and every
    candidate is scored under the same protocol.

  \item \textbf{Discovered algorithms that are architecturally new and generalize across sites.}
    We specify in advance what distinguishes an architectural change from a tuning change---whether
    the network itself differs from the starting paradigms, rather than carrying additional penalty
    terms or a different training schedule---and assess the selected candidates against it.

  \item \textbf{An interpretable account of what the discovered algorithms rely on.} We identify
    which connectome views and which regional pairs drive the prediction, and where these agree
    with established depression imaging findings and where they do not, which yields hypotheses
    that can be tested independently.
    %% RISK: 本条依赖演化结果。若最终胜出候选不可解释（例如闭式核方法），
    %% 必须撤下本条或改写为「对候选群体特征使用分布的分析」。定稿前按 §6.3 实际结果回改。

  \item \textbf{An evaluation protocol built to survive automated search.} The protocol uses
    site-disjoint partitions that keep model selection and reporting on mutually exclusive sets of
    scanners, together with a final test set that the search never reads and that is used once.
\end{itemize}

\paragraph{Roadmap.}
Section~\ref{sec:related} reviews reported classification results, the model families available for
connectome data, and evolutionary coding agents. Section~\ref{sec:agent} describes the agent.
Section~\ref{sec:application} specifies the data and the evaluation protocol, together with the
requirements that the difficulty of the task imposes on any such protocol.
Section~\ref{sec:results} reports results, and Section~\ref{sec:discussion} discusses limitations.
